% ============================================================
% 图论 C++ 知识点总结 —— 开卷考试参考
% 编译方式: xelatex 图论C++知识点总结.tex
% ============================================================
\documentclass[12pt,a4paper]{article}

\usepackage{geometry}
\geometry{left=2.0cm,right=2.0cm,top=2.0cm,bottom=2.0cm}

\usepackage{ctex}         % 中文支持
\usepackage{listings}     % 代码排版
\usepackage{xcolor}
\usepackage{hyperref}
\usepackage{amsmath,amssymb}
\usepackage{enumitem}
\usepackage{array}

% ===== 代码样式（与 study/CS/hyper/ 风格一致）=====
\lstset{
    language=C++,
    basicstyle=\small\ttfamily,
    keywordstyle=\bfseries\color{blue!75!black},
    commentstyle=\color{green!60!black},
    stringstyle=\color{red!70!black},
    numbers=left,
    numberstyle=\tiny\color{gray},
    numbersep=8pt,
    frame=none,
    breaklines=true,
    tabsize=4,
    showstringspaces=false,
    extendedchars=true,
    inputencoding=utf8,
    aboveskip=6pt,
    belowskip=6pt
}

% ===== 标题格式（参考 PDF：简单数字编号）=====
\usepackage{titlesec}
\titleformat{\section}{\Large\bfseries}{\thesection}{0.5em}{}
\titleformat{\subsection}{\large\bfseries}{\thesubsection}{0.5em}{}
\titleformat{\subsubsection}{\normalsize\bfseries}{\thesubsubsection}{0.5em}{}

\begin{document}

\begin{center}
    {\Huge\bfseries 图论 C++ 知识点总结}\\[0.3cm]
    {\large 数据结构考试复习}\\[0.2cm]
    {2026 年 6 月 18 日}
\end{center}
\vspace{0.3cm}

\tableofcontents
\newpage

% ============================================================
% 1  图的存储
% ============================================================
\section{图的存储}

\subsection{邻接矩阵（Adjacency Matrix）}
适用于稠密图，空间复杂度 $O(n^2)$。

\begin{lstlisting}[caption={邻接矩阵存储}]
const int N = 1005;
const int INF = 0x3f3f3f3f;

int graph[N][N];  // graph[u][v] = 边权

// 初始化
void initGraph(int n) {
    for (int i = 1; i <= n; ++i)
        for (int j = 1; j <= n; ++j)
            graph[i][j] = (i == j) ? 0 : INF;
}

// 添加有向边
void addEdge(int u, int v, int w) {
    graph[u][v] = w;
}

// 添加无向边
void addUndirectedEdge(int u, int v, int w) {
    graph[u][v] = w;
    graph[v][u] = w;
}
\end{lstlisting}

\subsection{邻接表（Adjacency List）}
适用于稀疏图，空间复杂度 $O(n + m)$。C++ 常用 \texttt{vector} 实现。

\begin{lstlisting}[caption={邻接表存储（vector 实现）}]
#include <vector>
using namespace std;

// 无权图：用 vector<vector<int>>
vector<vector<int>> adj;

void initGraph(int n) {
    adj.resize(n);
}

void addEdge(int u, int v) {
    adj[u].push_back(v);
    adj[v].push_back(u);  // 无向图
}

// 有权图
struct Edge {
    int to, w;
};
vector<vector<Edge>> adj2;

void initWeightedGraph(int n) {
    adj2.resize(n);
}

void addWeightedEdge(int u, int v, int w) {
    adj2[u].push_back({v, w});
    adj2[v].push_back({u, w});  // 无向图
}

// 遍历邻接点
for (int v : adj[u]) { /* ... */ }
for (Edge& e : adj2[u]) { /* e.to, e.w */ }
\end{lstlisting}

\subsection{结构体封装图类}

\begin{lstlisting}[caption={图类封装}]
struct ALGraph {
    vector<int> vertices[20];
    int vernum, arcnum;
};

void CreateGraph(ALGraph& G, int n, int e) {
    G.vernum = n;
    G.arcnum = e;
    for (int i = 0; i < e; i++) {
        int u, v;
        cin >> u >> v;
        G.vertices[u].push_back(v);
        G.vertices[v].push_back(u);  // 无向图
    }
}
\end{lstlisting}

\subsection{边集数组（用于 Kruskal 算法）}

\begin{lstlisting}[caption={边集存储}]
struct Edge {
    int u, v, w;
    bool operator<(const Edge& other) const {
        return w < other.w;  // 按权值排序
    }
};

vector<Edge> edges;
edges.push_back({u, v, w});
sort(edges.begin(), edges.end());
\end{lstlisting}

\newpage

% ============================================================
% 2  并查集
% ============================================================
\section{并查集（Union-Find / Disjoint Set Union）}
用于管理元素分组，支持查找和合并操作，近似 $O(\alpha(n))$。

\subsection{标准实现（路径压缩 + 按秩合并）}

\begin{lstlisting}[caption={并查集结构体}]
struct UnionFind {
    vector<int> parent;  // 父节点
    vector<int> rank;    // 秩（树高）

    UnionFind(int n) {
        parent.resize(n);
        rank.resize(n, 0);
        for (int i = 0; i < n; i++)
            parent[i] = i;
    }

    // 查找根 + 路径压缩
    int find(int x) {
        if (parent[x] != x)
            parent[x] = find(parent[x]);
        return parent[x];
    }

    // 按秩合并
    void unite(int x, int y) {
        int rx = find(x), ry = find(y);
        if (rx == ry) return;
        if (rank[rx] < rank[ry])
            parent[rx] = ry;
        else if (rank[rx] > rank[ry])
            parent[ry] = rx;
        else {
            parent[ry] = rx;
            rank[rx]++;
        }
    }

    bool same(int x, int y) {
        return find(x) == find(y);
    }
};
\end{lstlisting}

\subsection{简化实现（仅路径压缩）}

\begin{lstlisting}[caption={简单并查集}]
vector<int> parent;

void init(int n) {
    parent.resize(n + 1);
    for (int i = 1; i <= n; i++)
        parent[i] = i;
}

int find(int x) {
    if (parent[x] != x)
        parent[x] = find(parent[x]);
    return parent[x];
}

void unite(int x, int y) {
    int rx = find(x), ry = find(y);
    if (rx != ry) parent[ry] = rx;
}

bool same(int x, int y) {
    return find(x) == find(y);
}
\end{lstlisting}

\subsection{应用：统计连通分量}

\begin{lstlisting}[caption={并查集统计连通分量}]
// 不同的根节点数 = 连通分量数
int countComponents(int n) {
    init(n);
    for (每条边(u, v)) unite(u, v);
    int cnt = 0;
    for (int i = 1; i <= n; i++)
        if (find(i) == i) cnt++;
    return cnt;
}
\end{lstlisting}

\subsection{应用：分组处理}

\begin{lstlisting}[caption={并查集分组（如家族问题）}]
map<int, vector<int>> groups;
for (int i = 0; i < n; i++) {
    int root = uf.find(i);
    groups[root].push_back(i);
}
for (auto& entry : groups) {
    // entry.first = 根节点
    // entry.second = 该组所有成员
}
\end{lstlisting}

\newpage

% ============================================================
% 3  图的遍历
% ============================================================
\section{图的遍历}

\subsection{深度优先搜索（DFS）}

\begin{lstlisting}[caption={DFS 递归 / 非递归}]
#include <vector>
#include <stack>
using namespace std;

// 递归
void dfs(int u, vector<vector<int>>& adj, vector<bool>& visited) {
    visited[u] = true;
    cout << u << " ";
    for (int v : adj[u])
        if (!visited[v]) dfs(v, adj, visited);
}

// 非递归（栈）
void dfsStack(int start, vector<vector<int>>& adj) {
    int n = adj.size();
    vector<bool> vis(n, false);
    stack<int> st;
    st.push(start);
    while (!st.empty()) {
        int u = st.top(); st.pop();
        if (vis[u]) continue;
        vis[u] = true;
        cout << u << " ";
        for (int i = adj[u].size() - 1; i >= 0; i--) {
            int v = adj[u][i];
            if (!vis[v]) st.push(v);
        }
    }
}
\end{lstlisting}

\subsection{广度优先搜索（BFS）}

\begin{lstlisting}[caption={BFS 实现（含非连通图处理）}]
#include <queue>
#include <vector>
#include <iostream>
using namespace std;

void bfs(int start, vector<vector<int>>& adj,
         vector<bool>& visited, vector<int>& traversal) {
    queue<int> q;
    q.push(start);
    visited[start] = true;
    while (!q.empty()) {
        int u = q.front(); q.pop();
        traversal.push_back(u);
        for (int v : adj[u]) {
            if (!visited[v]) {
                visited[v] = true;
                q.push(v);
            }
        }
    }
}

// 非连通图：每个未访问点为起点做 BFS
void graphBFS(int n, vector<vector<int>>& adj) {
    vector<bool> visited(n, false);
    vector<int> traversal;
    int components = 0;
    for (int i = 0; i < n; i++) {
        if (!visited[i]) {
            bfs(i, adj, visited, traversal);
            components++;
        }
    }
    // traversal: BFS 序列
    // components: 连通分量数
}
\end{lstlisting}

\subsection{BFS 求最短路径（无权图）}

\begin{lstlisting}[caption={BFS 求最短路径步数}]
vector<int> bfsShortestPath(int start, int n, vector<vector<int>>& adj) {
    vector<int> dist(n, -1);  // -1 表示不可达
    queue<int> q;
    dist[start] = 0;
    q.push(start);
    while (!q.empty()) {
        int u = q.front(); q.pop();
        for (int v : adj[u]) {
            if (dist[v] == -1) {
                dist[v] = dist[u] + 1;
                q.push(v);
            }
        }
    }
    return dist;
}
\end{lstlisting}

\subsection{迷宫 BFS（二维网格）}

\begin{lstlisting}[caption={二维迷宫 / 网格 BFS}]
struct Point { int x, y, step; };

int bfsMaze(vector<vector<int>>& maze, int sx, int sy,
            int ex, int ey) {
    int n = maze.size(), m = maze[0].size();
    int dx[4] = {-1, 1, 0, 0};
    int dy[4] = {0, 0, -1, 1};

    queue<Point> q;
    vector<vector<bool>> vis(n, vector<bool>(m, false));
    q.push({sx, sy, 0});
    vis[sx][sy] = true;

    while (!q.empty()) {
        Point cur = q.front(); q.pop();
        if (cur.x == ex && cur.y == ey) return cur.step;

        for (int i = 0; i < 4; i++) {
            int nx = cur.x + dx[i], ny = cur.y + dy[i];
            if (nx < 0 || nx >= n || ny < 0 || ny >= m) continue;
            if (maze[nx][ny] == 1 || vis[nx][ny]) continue;
            vis[nx][ny] = true;
            q.push({nx, ny, cur.step + 1});
        }
    }
    return -1;  // 不可达
}
\end{lstlisting}

\subsection{三维 BFS}

\begin{lstlisting}[caption={三维 BFS（6 方向）}]
struct Point3D { int x, y, z; };
int dx[6] = {1, -1, 0, 0, 0, 0};
int dy[6] = {0, 0, 1, -1, 0, 0};
int dz[6] = {0, 0, 0, 0, 1, -1};

int bfs3D(vector<vector<vector<int>>>& mat,
          vector<vector<vector<bool>>>& vis,
          int M, int N, int L, int sx, int sy, int sz) {
    queue<Point3D> q;
    q.push({sx, sy, sz});
    vis[sx][sy][sz] = true;
    int vol = 1;
    while (!q.empty()) {
        Point3D p = q.front(); q.pop();
        for (int i = 0; i < 6; i++) {
            int nx = p.x + dx[i], ny = p.y + dy[i], nz = p.z + dz[i];
            if (nx < 0 || nx >= M || ny < 0 || ny >= N || nz < 0 || nz >= L)
                continue;
            if (mat[nx][ny][nz] == 1 && !vis[nx][ny][nz]) {
                vis[nx][ny][nz] = true;
                q.push({nx, ny, nz});
                vol++;
            }
        }
    }
    return vol;
}
\end{lstlisting}

\newpage

% ============================================================
% 4  拓扑排序
% ============================================================
\section{拓扑排序（Topological Sort）}
有向无环图（DAG）的顶点线性排序，每条有向边的起点都在终点之前。

\subsection{Kahn 算法（BFS 队列实现）}

\begin{lstlisting}[caption={拓扑排序（Kahn 算法）}]
#include <iostream>
#include <vector>
#include <queue>
using namespace std;

vector<int> topologicalSort(int n, vector<vector<int>>& adj) {
    vector<int> indeg(n + 1, 0);
    for (int u = 1; u <= n; u++)
        for (int v : adj[u]) indeg[v]++;

    queue<int> q;
    for (int i = 1; i <= n; i++)
        if (indeg[i] == 0) q.push(i);

    vector<int> result;
    while (!q.empty()) {
        int u = q.front(); q.pop();
        result.push_back(u);
        for (int v : adj[u]) {
            indeg[v]--;
            if (indeg[v] == 0) q.push(v);
        }
    }

    // 结果长度 != n 表示有环
    return result.size() == n ? result : vector<int>{};
}
\end{lstlisting}

\subsection{栈实现}

\begin{lstlisting}[caption={拓扑排序（栈实现）}]
vector<int> topoSortByStack(int n, vector<vector<int>>& adj) {
    vector<int> indeg(n + 1, 0);
    for (int u = 1; u <= n; u++)
        for (int v : adj[u]) indeg[v]++;

    stack<int> st;
    for (int i = 1; i <= n; i++)
        if (indeg[i] == 0) st.push(i);

    int count = 0;
    vector<int> result;
    while (!st.empty()) {
        int u = st.top(); st.pop();
        result.push_back(u);
        count++;
        for (int v : adj[u]) {
            indeg[v]--;
            if (indeg[v] == 0) st.push(v);
        }
    }
    if (count != n) return {};  // 有环
    return result;
}
\end{lstlisting}

\newpage

% ============================================================
% 5  最短路径
% ============================================================
\section{最短路径}

\subsection{Dijkstra 算法（单源最短路，非负权）}

\subsubsection{朴素版 $O(n^2)$}

\begin{lstlisting}[caption={Dijkstra 朴素版}]
#include <iostream>
#include <vector>
#include <climits>
using namespace std;

const int INF = INT_MAX / 2;

vector<int> dijkstra(int n, int start, vector<vector<int>>& graph) {
    vector<int> dist(n + 1, INF);
    vector<bool> visited(n + 1, false);
    dist[start] = 0;

    for (int i = 1; i <= n; i++) {
        int u = -1, minDist = INF;
        for (int j = 1; j <= n; j++)
            if (!visited[j] && dist[j] < minDist) {
                minDist = dist[j];
                u = j;
            }
        if (u == -1) break;
        visited[u] = true;

        for (int v = 1; v <= n; v++)
            if (!visited[v] && graph[u][v] != INF
                && dist[u] + graph[u][v] < dist[v])
                dist[v] = dist[u] + graph[u][v];
    }
    return dist;
}
\end{lstlisting}

\subsubsection{堆优化版 $O(m\log n)$}

\begin{lstlisting}[caption={Dijkstra 堆优化}]
#include <queue>
#include <vector>
using namespace std;

const int INF = 0x3f3f3f3f;
struct Edge { int to, w; };
struct Node {
    int id, dist;
    bool operator>(const Node& o) const {
        return dist > o.dist;
    }
};

vector<vector<Edge>> adj;

vector<int> dijkstraHeap(int n, int start) {
    adj.resize(n + 1);
    vector<int> dist(n + 1, INF);
    vector<bool> vis(n + 1, false);
    priority_queue<Node, vector<Node>, greater<Node>> pq;

    dist[start] = 0;
    pq.push({start, 0});

    while (!pq.empty()) {
        Node cur = pq.top(); pq.pop();
        int u = cur.id;
        if (vis[u]) continue;
        vis[u] = true;
        for (Edge& e : adj[u]) {
            int v = e.to, w = e.w;
            if (dist[u] + w < dist[v]) {
                dist[v] = dist[u] + w;
                pq.push({v, dist[v]});
            }
        }
    }
    return dist;
}
\end{lstlisting}

\subsection{Floyd-Warshall 算法（多源最短路）}
$O(n^3)$，适用于 $n \leq 500$。

\begin{lstlisting}[caption={Floyd 算法}]
const int INF = 0x3f3f3f3f;
int dist[505][505];

void floyd(int n) {
    for (int k = 1; k <= n; k++)
        for (int i = 1; i <= n; i++)
            for (int j = 1; j <= n; j++)
                if (dist[i][k] != INF && dist[k][j] != INF)
                    dist[i][j] = min(dist[i][j],
                                     dist[i][k] + dist[k][j]);
}
\end{lstlisting}

\subsection{Bellman-Ford 算法（含负权）}

\begin{lstlisting}[caption={Bellman-Ford 算法}]
struct Edge { int u, v, w; };
vector<Edge> edges;

bool bellmanFord(int n, int start, vector<int>& dist) {
    dist.assign(n + 1, INF);
    dist[start] = 0;
    for (int i = 1; i < n; i++) {
        bool updated = false;
        for (Edge& e : edges)
            if (dist[e.u] != INF && dist[e.u] + e.w < dist[e.v]) {
                dist[e.v] = dist[e.u] + e.w;
                updated = true;
            }
        if (!updated) break;
    }
    // 第 n 轮检测负环
    for (Edge& e : edges)
        if (dist[e.u] != INF && dist[e.u] + e.w < dist[e.v])
            return false;  // 有负环
    return true;
}
\end{lstlisting}

\newpage

% ============================================================
% 6  最小生成树
% ============================================================
\section{最小生成树（Minimum Spanning Tree）}

\subsection{Kruskal 算法（加边法，$O(m\log m)$）}
基于并查集，按边权从小到大选取。

\begin{lstlisting}[caption={Kruskal 算法}]
#include <algorithm>
#include <vector>
using namespace std;

struct Edge {
    int u, v, w;
    bool operator<(const Edge& o) const { return w < o.w; }
};

// 并查集（见第 2 节）
vector<int> parent;
int find(int x) { ... }
void unite(int x, int y) { ... }

int kruskal(int n, vector<Edge>& edges) {
    parent.resize(n + 1);
    for (int i = 1; i <= n; i++) parent[i] = i;

    sort(edges.begin(), edges.end());
    int mstWeight = 0, edgeCount = 0;

    for (Edge& e : edges) {
        if (find(e.u) != find(e.v)) {
            unite(e.u, e.v);
            mstWeight += e.w;
            edgeCount++;
            if (edgeCount == n - 1) break;
        }
    }
    return edgeCount == n - 1 ? mstWeight : -1;  // -1 表示不连通
}
\end{lstlisting}

\subsection{Prim 算法（加点法，$O(n^2)$）}

\begin{lstlisting}[caption={Prim 朴素版}]
const int INF = 0x3f3f3f3f;
int graph[505][505];

int prim(int n) {
    vector<int> dist(n + 1, INF);
    vector<bool> visited(n + 1, false);
    dist[1] = 0;
    int mstWeight = 0;

    for (int i = 1; i <= n; i++) {
        int u = -1, minDist = INF;
        for (int j = 1; j <= n; j++)
            if (!visited[j] && dist[j] < minDist) {
                minDist = dist[j];
                u = j;
            }
        if (u == -1) return -1;  // 不连通
        visited[u] = true;
        mstWeight += dist[u];

        for (int v = 1; v <= n; v++)
            if (!visited[v] && graph[u][v] < dist[v])
                dist[v] = graph[u][v];
    }
    return mstWeight;
}
\end{lstlisting}

\newpage

% ============================================================
% 7  关键路径
% ============================================================
\section{关键路径（AOE 网）}

\begin{lstlisting}[caption={拓扑排序求关键路径}]
struct Edge { int to, w; };

// 正推求 ve（最早发生时间）
vector<int> topoOrder(int n, vector<vector<Edge>>& adj,
                       vector<int>& indeg, vector<int>& ve) {
    queue<int> q;
    for (int i = 1; i <= n; i++)
        if (indeg[i] == 0) q.push(i);

    vector<int> topo;
    while (!q.empty()) {
        int u = q.front(); q.pop();
        topo.push_back(u);
        for (Edge& e : adj[u]) {
            indeg[e.to]--;
            if (indeg[e.to] == 0) q.push(e.to);
            ve[e.to] = max(ve[e.to], ve[u] + e.w);
        }
    }
    return topo;
}

// 逆推求 vl（最晚发生时间）
void calcVl(int n, vector<int>& topo, vector<vector<Edge>>& adj,
            int totalTime, vector<int>& vl) {
    fill(vl.begin(), vl.begin() + n + 1, totalTime);
    for (int i = n - 1; i >= 0; i--) {
        int u = topo[i];
        for (Edge& e : adj[u])
            vl[u] = min(vl[u], vl[e.to] - e.w);
    }
}

// 关键活动：ve[u] == vl[u] 且 ve[u] + w == vl[v]
\end{lstlisting}

\newpage

% ============================================================
% 8  常见技巧与模板
% ============================================================
\section{常见技巧与模板}

\subsection{快速 I/O}

\begin{lstlisting}[caption={C++ 加速 I/O}]
ios_base::sync_with_stdio(false);
cin.tie(nullptr);
// 注意：不可混用 scanf/printf 和 cin/cout
\end{lstlisting}

\subsection{INF 常量}

\begin{lstlisting}
const int INF1 = 0x3f3f3f3f;    // ≈ 1.06e9，两个相加不溢出
const int INF2 = INT_MAX / 2;   // 更保守
\end{lstlisting}

\subsection{重边处理}

\begin{lstlisting}
if (w < graph[u][v]) {    // 保留最小权值
    graph[u][v] = w;
    graph[v][u] = w;      // 无向图
}
\end{lstlisting}

\subsection{邻接表排序}

\begin{lstlisting}
for (int i = 0; i < n; i++)
    sort(adj[i].begin(), adj[i].end());
\end{lstlisting}

\subsection{DFS 检测有向环}

\begin{lstlisting}[caption={三色标记法}]
bool hasCycle(int u, vector<vector<int>>& adj, vector<int>& color) {
    color[u] = 1;  // 正在访问
    for (int v : adj[u]) {
        if (color[v] == 1) return true;       // 后向边 → 有环
        if (color[v] == 0 && hasCycle(v, adj, color)) return true;
    }
    color[u] = 2;  // 已访问完
    return false;
}

// 调用：
// vector<int> color(n, 0);
// for (int i = 0; i < n; i++)
//     if (color[i] == 0 && hasCycle(i, adj, color)) ...
\end{lstlisting}

\subsection{索引排序}

\begin{lstlisting}[caption={按值排序返回原下标}]
vector<int> isort(const vector<int>& arr, int n) {
    vector<int> ia(n);
    for (int i = 0; i < n; i++) ia[i] = i;
    sort(ia.begin(), ia.end(), [&](int a, int b) {
        return arr[a] < arr[b];
    });
    return ia;
}
// ia[0] = arr 中最小值的下标
\end{lstlisting}

\subsection{优先队列（堆）}

\begin{lstlisting}[caption={priority\_queue 用法}]
// 最大堆（默认）
priority_queue<int> pq;

// 最小堆
priority_queue<int, vector<int>, greater<int>> pq_min;

// 自定义类型（小根堆）
struct Node {
    int dist, id;
    bool operator>(const Node& o) const {
        return dist > o.dist;
    }
};
priority_queue<Node, vector<Node>, greater<Node>> pq;
\end{lstlisting}

\newpage

% ============================================================
% 附录：算法速查表
% ============================================================
\section*{附录：算法速查表}
\addcontentsline{toc}{section}{附录：算法速查表}

\begin{center}
\begin{tabular}{|c|c|c|}
\hline
\textbf{算法} & \textbf{时间复杂度} & \textbf{核心数据结构} \\
\hline
并查集 Find & $O(\alpha(n))$ & parent[] 数组 \\
\hline
并查集 Unite & $O(\alpha(n))$ & rank[] 数组 \\
\hline
DFS & $O(n + m)$ & 递归 / stack \\
\hline
BFS & $O(n + m)$ & queue \\
\hline
拓扑排序 & $O(n + m)$ & indeg[] + queue/stack \\
\hline
Dijkstra 朴素 & $O(n^2)$ & dist[] + visited[] \\
\hline
Dijkstra 堆优化 & $O(m\log n)$ & priority\_queue \\
\hline
Floyd & $O(n^3)$ & dist[][] \\
\hline
Bellman-Ford & $O(nm)$ & Edge[] \\
\hline
Kruskal & $O(m\log m)$ & Union-Find \\
\hline
Prim 朴素 & $O(n^2)$ & dist[] + visited[] \\
\hline
\end{tabular}
\end{center}

\vfill
\begin{flushright}
\small —— 祝考试顺利！——
\end{flushright}

\end{document}
